\chapter{Metodología}
Modelos que vamos a probar, preprocesado de datos (normalización, etc)
\section{Extracción de Keypoints}
Como se mencionó, la extracción de keypoints es el proceso mediante el cual se obtiene una representación estructurada del cuerpo humano a partir de imágenes o vídeo, en forma de coordenadas de puntos anatómicos predefinidos, tales como articulaciones, extremidades o rasgos faciales. Esta representación abstrae la información visual relevante para el movimiento y la postura, descartando información irrelevante como el fondo, la iluminación o la apariencia superficial, lo que la convierte en una representación compacta, generalizable y respetuosa con la privacidad, tal y como se argumentó en la sección anterior.

En los últimos años, el avance conjunto de la visión por computador y el aprendizaje profundo ha dado lugar a herramientas de estimación de pose de gran precisión y eficiencia. Entre las más destacadas se encuentran \textbf{MediaPipe Holistic}, desarrollada por Google y ampliamente adoptada en la literatura de lengua de signos por su capacidad para detectar simultáneamente cuerpo, manos y rostro en tiempo real con un \textit{bajo costo en recursos computacionales}; \textbf{OpenPose}, uno de los primeros sistemas de estimación multi-persona de referencia, si bien \textit{ha sido ampliamente utilizado como referencia} en trabajos previos de estimación de pose, presenta limitaciones en términos de eficiencia y robustez frente a enfoques más recientes; aunque métodos como \textbf{HRNet} ofrecen una \textit{elevada precisión en la localización de keypoints en imágenes estáticas}, no incorporan de forma explícita mecanismos de seguimiento temporal, lo que puede generar variabilidad entre frames en secuencias de vídeo. Finalmente, \textbf{AlphaPose}, un sistema de estimación de pose de alta precisión que destaca especialmente en escenarios con múltiple personas, \textit{condiciones de oclusión y videos de baja calidad}, que es espeicalmete nuestro caso. Además, este último, a diferencia de \textbf{HRNet}, está optimizado para ser relativamente rápido, permitiendo una extracción más cómoda.




%\subsection{OpenPose}
\subsection{Alphapose}
AlphaPose ha demostrado obtener el menor error medio en detección de keypoints frente a MediaPipe y OpenPose incluso bajo condiciones de degradación de imagen \cite{info16110970}, lo que refuerza su idoneidad para entornos de grabación no controlados, habituales en datasets de lengua de signos, especialmente en el dataset con el que trabajaremos.

AlphaPose ofrece la capacidad de cambiar el modelo utilizado para la representación del esqueleto, existiendo principalmente dos conjuntos de datos de pose de cuerpo completo:
\begin{itemize}
  \item \textbf{COCO-WholeBody}, una ampliación del conjunto de datos \textbf{MS COCO} enfocada en el cuerpo humano completo. Incluye anotaciones detalladas de puntos clave del cuerpo, manos, cara y pies, lo que permite realizar estimaciones de pose más completas y precisas en tareas de análisis del movimiento humano.
  \item \textbf{HALPE-136}, un conjunto de datos más reciente que proporciona una anotación aún más densa (136 puntos clave), con especial énfasis en manos y rostro. Este mayor nivel de detalle puede resultar beneficioso en tareas que requieren capturar movimientos finos, como el reconocimiento de lengua de signos, aunque a costa de una mayor complejidad computacional y posible sensibilidad al ruido en entornos no controlados.
\end{itemize}
\begin{figure}[htbp]
  \centering
  % --- BLOQUE IZQUIERDO: La figura alta ---
  \begin{minipage}[b]{0.48\textwidth}
    \centering
    \includegraphics[width=\textwidth]{human_model_halpe_136.jpg}

  \end{minipage}
  \hfill
  % --- BLOQUE DERECHO: Las dos figuras cuadradas ---
  \begin{minipage}[b]{0.48\textwidth}
    \centering
    % Primera figura cuadrada
    \includegraphics[width=0.8\textwidth]{hand_halpe_136.jpg}

    \vspace{0.5cm}

    % Segunda figura cuadrada
    \includegraphics[width=0.8\textwidth]{face_halpe_136.jpg}

  \end{minipage}
  \caption{Keypoints en el Modelo Halpe-136}
\end{figure}

\lstinputlisting[language=json]{code/keypoints-format.txt}


\section{Seq2Seq}

\section{Normalización de los Datos}

\section{Multi Stream Keypoint Attention}
En esta sección se presentará la arquitectura principal del modelo \textit{Multi-stream Keypoint Attention} (MSKA) detallando el razonamiento detrás de la misma. Es importante notar que este modelo \textbf{no} es \textit{end-to-end}, hace uso de las glosas como representación intermedia de la traducción, asegurando que el modelo realmente aprende el significado detrás de los \textit{keypoints}.



